% Companion perspective paper to paper/goo/main.tex: the other ways a
% civilisation can die of its own self-replicating technology, organised by
% the three properties that govern reach and visibility.  Numbers shared with
% the main paper are read from ../goo/numbers.tex.
\documentclass[11pt]{article}
\usepackage[margin=1in]{geometry}
\usepackage{amsmath,amssymb,graphicx,natbib,booktabs,hyperref,xspace,xcolor}
\newcommand{\um}{\ensuremath{\,\mu\mathrm{m}}\xspace}
\newcommand{\kms}{\ensuremath{\,\mathrm{km\,s^{-1}}}\xspace}
\graphicspath{{../../results/goo/figures/}}
\bibliographystyle{plainnat}
% generated by `python -m seti.goo.run --stage numbers`; do not edit
\newcommand{\SeedMassKg}{1.0\times10^{-15}\xspace}
\newcommand{\SeedRadiusUm}{0.5\xspace}
\newcommand{\ERepJkg}{1\times10^{7}\xspace}
\newcommand{\BioDoublings}{100\xspace}
\newcommand{\CrustDoublings}{124\xspace}
\newcommand{\BulkDoublings}{132\xspace}
\newcommand{\BioExpTime}{27.7\,h\xspace}
\newcommand{\BioRadTime}{1.8\,h\xspace}
\newcommand{\BioSolarTime}{25.0\,h\xspace}
\newcommand{\CrustRadTime}{5\,kyr\xspace}
\newcommand{\CrustSolarTime}{62\,kyr\xspace}
\newcommand{\CrustGovTime}{5\,kyr\xspace}
\newcommand{\BulkRadTime}{1.1\,Myr\xspace}
\newcommand{\BulkSolarTime}{15.5\,Myr\xspace}
\newcommand{\BulkGovTime}{1.1\,Myr\xspace}
\newcommand{\BulkRadTimeHot}{4\,kyr\xspace}
\newcommand{\CrustRadTimeHot}{16\,yr\xspace}
\newcommand{\RadLimitFiveHundred}{1.7\times10^{18}\xspace}
\newcommand{\RadLimitTwoThousand}{4.6\times10^{20}\xspace}
\newcommand{\AbsorbedSolar}{1.2\times10^{17}\xspace}
\newcommand{\CrustFissionJ}{8\times10^{30}\xspace}
\newcommand{\OceanDeuteriumJ}{7\times10^{31}\xspace}
\newcommand{\EarthDisassemblyLsun}{7\,d\xspace}
\newcommand{\EarthDisassemblySolar}{58.2\,Myr\xspace}
\newcommand{\EarthDisassemblyRad}{15\,kyr\xspace}
\newcommand{\BioExpTimeTauHundred}{2.8\,h\xspace}
\newcommand{\CrustGovTimeTauHundred}{5\,kyr\xspace}
\newcommand{\BioExpTimeTauTenK}{12\,d\xspace}
\newcommand{\CrustGovTimeTauTenK}{5\,kyr\xspace}
\newcommand{\BioExpTimeTauYear}{101\,yr\xspace}
\newcommand{\CrustGovTimeTauYear}{5\,kyr\xspace}
\newcommand{\BeltMassKg}{2.4\times10^{21}\xspace}
\newcommand{\BeltConvTime}{14\,yr\xspace}
\newcommand{\KuiperConvTime}{1\,kyr\xspace}
\newcommand{\BeltTempK}{169\xspace}
\newcommand{\KuiperTempK}{44\xspace}
\newcommand{\BeltCoverFilm}{1\xspace}
\newcommand{\BeltCoverSheet}{0.0012\xspace}
\newcommand{\BeltCoverSlab}{1.2\times10^{-6}\xspace}
\newcommand{\KuiperCoverSheet}{0.00027\xspace}
\newcommand{\TerrConvTime}{8\,kyr\xspace}
\newcommand{\TerrConvTimeFullL}{51.6\,h\xspace}
\newcommand{\GiantConvTime}{143\,kyr\xspace}
\newcommand{\GiantConvTimeFullL}{21\,yr\xspace}
\newcommand{\BeltResidueMicron}{4\,yr\xspace}
\newcommand{\BeltResidueMm}{10\,kyr\xspace}
\newcommand{\BeltResidueMetre}{10.1\,Myr\xspace}
\newcommand{\BeltCollTime}{5\,yr\xspace}
\newcommand{\BeltResidueMmCover}{4.4\times10^{-4}\xspace}
\newcommand{\BeltResidueMetreCover}{4.4\times10^{-7}\xspace}
\newcommand{\MaxVisibleResidueMilli}{4\,kyr\xspace}
\newcommand{\MaxVisibleResidueCenti}{444\,yr\xspace}
\newcommand{\MaxVisibleResidueDeci}{44\,yr\xspace}
\newcommand{\KuiperMaxVisibleResidueCenti}{25\,kyr\xspace}
\newcommand{\HephBoundTenK}{6\xspace}
\newcommand{\HephBoundMyr}{0.06\xspace}
\newcommand{\HephBoundTenGyr}{6.0\times10^{-6}\xspace}
\newcommand{\GhatFractionBound}{3.0\times10^{-5}\xspace}
\newcommand{\BlowoutRadiusUm}{0.57\xspace}
\newcommand{\BetaMaxReach}{1.38\xspace}
\newcommand{\WindowLoUm}{0.21\xspace}
\newcommand{\WindowHiUm}{0.57\xspace}
\newcommand{\VinfDevice}{7\xspace}
\newcommand{\VinfBetaOne}{18\xspace}
\newcommand{\AnnulusFillGyr}{1.2\xspace}
\newcommand{\FocusingFactor}{1.2\xspace}
\newcommand{\BlownFractionColl}{3.0\times10^{-4}\xspace}
\newcommand{\BeltDevicesReleased}{2.3\times10^{36}\xspace}
\newcommand{\BeltArrivalsUnsurvived}{2e+02\xspace}
\newcommand{\BeltCumulativeMyr}{0\xspace}
\newcommand{\BeltCumulativeTenMyr}{1.3\times10^{-42}\xspace}
\newcommand{\BeltCumulativeHundredMyr}{1.5\times10^{5}\xspace}
\newcommand{\BeltCumulativeGyr}{5.9\times10^{10}\xspace}
\newcommand{\BioArrivalsUnsurvived}{8.1\times10^{-5}\xspace}
\newcommand{\BioCumulativeGyr}{2.5\times10^{4}\xspace}
\newcommand{\BioCumulativeHundredMyr}{0.062\xspace}
\newcommand{\FillSlowSlow}{504.2\,Myr\xspace}
\newcommand{\FillMidMid}{6.4\,Myr\xspace}
\newcommand{\FillFastFast}{639\,kyr\xspace}
\newcommand{\FillVoyagerTen}{489.4\,Myr\xspace}
\newcommand{\FrontMidMid}{2295\xspace}
\newcommand{\EventHorizonGly}{16.6\xspace}
\newcommand{\ReachTenthGly}{1.7\xspace}
\newcommand{\GalaxiesTenth}{5\times10^{6}\xspace}
\newcommand{\ReachHalfGly}{8.3\xspace}
\newcommand{\GalaxiesHalf}{7\times10^{8}\xspace}
\newcommand{\ReachEightTenthsGly}{13.2\xspace}
\newcommand{\GalaxiesEightTenths}{3\times10^{9}\xspace}
\newcommand{\ReachNinetyNineGly}{16.4\xspace}
\newcommand{\GalaxiesNinetyNine}{5\times10^{9}\xspace}
\newcommand{\ReachLightGly}{16.6\xspace}
\newcommand{\GalaxiesLight}{5\times10^{9}\xspace}
\newcommand{\WindowGyr}{5\xspace}
\newcommand{\NGalObserved}{1\times10^{5}\xspace}
\newcommand{\PnaiveLamOne}{0.13\xspace}
\newcommand{\PshadowLamOne}{0.25\xspace}
\newcommand{\PextPersistLamOne}{3.6\times10^{-6}\xspace}
\newcommand{\PextDeadLamOne}{0.02\xspace}
\newcommand{\PpassiveNaiveLamOne}{0.16\xspace}
\newcommand{\PnaiveLamHundred}{2.3\times10^{-3}\xspace}
\newcommand{\PshadowLamHundred}{0.25\xspace}
\newcommand{\PextPersistLamHundred}{5.1\times10^{-8}\xspace}
\newcommand{\PextDeadLamHundred}{2.6\times10^{-4}\xspace}
\newcommand{\PpassiveNaiveLamHundred}{3.2\times10^{-3}\xspace}
\newcommand{\PnaiveLamTenK}{3.1\times10^{-5}\xspace}
\newcommand{\PshadowLamTenK}{0.25\xspace}
\newcommand{\PextPersistLamTenK}{7.9\times10^{-10}\xspace}
\newcommand{\PextDeadLamTenK}{3.6\times10^{-6}\xspace}
\newcommand{\PpassiveNaiveLamTenK}{4.3\times10^{-5}\xspace}
\newcommand{\LbOwnLamOne}{3\xspace}
\newcommand{\LbOwnPassiveLamOne}{4.3\xspace}
\newcommand{\LbExtPersistLamOne}{3.0\times10^{-5}\xspace}
\newcommand{\LbExtDeadLamOne}{0.3\xspace}
\newcommand{\LbOwnLamHundred}{0.03\xspace}
\newcommand{\LbOwnPassiveLamHundred}{0.043\xspace}
\newcommand{\LbExtPersistLamHundred}{3.0\times10^{-7}\xspace}
\newcommand{\LbExtDeadLamHundred}{3.0\times10^{-3}\xspace}
\newcommand{\LbOwnLamTenK}{3.0\times10^{-4}\xspace}
\newcommand{\LbOwnPassiveLamTenK}{4.3\times10^{-4}\xspace}
\newcommand{\LbExtPersistLamTenK}{3.0\times10^{-9}\xspace}
\newcommand{\LbExtDeadLamTenK}{3.0\times10^{-5}\xspace}
\newcommand{\PmargNaive}{0.031\xspace}
\newcommand{\PmargShadow}{0.25\xspace}
\newcommand{\PmargExtPersist}{8.3\times10^{-6}\xspace}
\newcommand{\PmargExtDead}{0.011\xspace}
\newcommand{\BranchRatioTypical}{0.990\xspace}
\newcommand{\BranchRatioHalf}{0.50\xspace}
\newcommand{\BranchRatioTenth}{0.90\xspace}


\title{Dying of one's own replicators: dispersal, reach and visibility for the
self-replicating technologies with which a civilisation can end itself}
\author{[Author list]}
\date{\today}

\begin{document}
\maketitle

\begin{abstract}
Molecular assemblers, engineered organisms, gene drives, self-reproducing
machines and space industry, network worms and self-replicating software
agents are all replicators, and every one of them has been proposed as a way
a technological civilisation might accidentally end itself.  They are usually
treated one at a time and by different communities.  We treat them together,
by the three properties that decide what an escaped replicator can do:
whether its growth has an external limit, how it disperses, and whether it
keeps metabolising where an outside observer can see.  The first sets whether
it is a catastrophe; the second sets its reach; the third sets whether the
universe keeps a record.  We show that accidental replicators inherit the
dispersal of the environment they escape into---the biosphere, the network,
the planetary system---and that for a civilisation at our stage every
plausible accident is bounded by its planet.  Only the replicator a
civilisation builds on purpose to travel between stars has a reach that
astronomy can see, and the observational record (compiled in a companion
paper) constrains that class and no other: fewer than one civilisation in
$10^{6}$--$10^{7}$ can have released a persistent galaxy-spanning replicator if
civilisations are common, while nothing in the sky bears on the planet-bound
classes at all.  We draw the consequences for how the classes should be
prioritised, for what ``containment'' means at each scale, and for the
peculiar epistemic position of a species trying to estimate a risk whose
realisation would leave no trace anyone could observe.
\end{abstract}

\section{Introduction}

The idea that a civilisation could be destroyed by something it built that
copies itself is older than any of the technologies now proposed to do it.
Von Neumann's universal constructor \citep{FreitasMerkle2004} made
self-replication an engineering concept in the 1940s; the Morris worm of 1988
\citep{Spafford1989} made it an incident; Drexler's ``grey goo''
\citep{Drexler1986} made it an extinction scenario, and Joy's essay
\citep{Joy2000} made it a public one.  The existential-risk literature since
\citet{Bostrom2002} has listed self-replicators under several headings---
nanotechnology, engineered pandemics, misaligned artificial
intelligence---and has assigned each a probability by elicitation
\citep{Sandberg2008,Ord2020}, while the technical communities that could
build them have argued, class by class, about whether the danger is real
\citep{PhoenixDrexler2004,Adamala2024,Noble2018,Pan2024}.

This paper is a perspective, not a survey.  Its claim is that the classes
share a structure that is more informative than any of them alone, and that
the structure has three parts.  First, a replicator is a catastrophe only if
its growth has no external limit; most replicators most of the time are
limited by feedstock, heat, competition or copying error, and the question
``can it escape?'' is prior to every other.  Second, an escaped replicator's
\emph{reach} is set by how it disperses, and it inherits that dispersal from
the environment it escapes into: a pathogen's reach is its host species', a
worm's is its network's, a molecular assembler's is its planet's unless it is
in space.  Third, a replicator's \emph{visibility} from outside---whether the
universe keeps a record of the catastrophe---is set by whether it keeps
metabolising where an outside observer can see it, which for astronomy means
re-radiating starlight around a star or across a galaxy.  The companion paper
\citep{GooPaper} works the molecular case through all three properties at
every scale from planet to cosmos and compiles the observational record; here
we apply the same three questions to every class and draw the comparative
conclusions.

\section{The three properties}

\paragraph{Limit.}  Every replicator population obeys some form of
$\dot N = rN\,(1 - N/K) - \mu N$: growth at rate $r$, a carrying capacity $K$
set by feedstock or heat, and a loss $\mu$ from decay, predation, competition
or copying error.  The error catastrophe \citep{Eigen1971} is the general
statement that a finite-fidelity copier of genome length $L$ and per-symbol
error $\epsilon$ loses its identity once $\epsilon L\gtrsim\ln s$ for
selective advantage $s$; biology sits just under that threshold and any
engineered replicator must be placed there deliberately.  A replicator is a
catastrophe only in the regime where $K$ is the whole environment and $\mu$ is
small, which is why designed replicators come with limiters
\citep{Ellery2022b} and why the accident is the failure of a limiter rather
than the design of a monster.

\paragraph{Dispersal.}  Once loose, the population's reach is
$\min(\text{substrate extent},\ \text{dispersal extent})$.  The companion paper
shows for the molecular case that dispersal class, not replication rate,
decides everything: a planet's crust is a \CrustGovTime project and its bulk a
\BulkGovTime one at any doubling time because the surface cannot shed the
heat faster, a planetary system's small bodies are transport-limited to
\BeltConvTime--\KuiperConvTime, and the Galaxy is reached only by a
replicator that is either smaller than \BlowoutRadiusUm\um and loose in space
(blown out on starlight, crossing the Galaxy as a chain of seedings in
\SurvGrainTGalInf\ if it stays functional for longer than the $\sim10^{5}$--$10^{6}$\,yr
hop to the next star) or in possession of a starship (\FillFastFast--%
\FillSlowSlow to fill the disc).  The same logic applies to every class.

\paragraph{Visibility.}  A catastrophe leaves an astronomical record only if
the replicator keeps converting energy where a survey can see: a live swarm
re-radiates starlight and is a Dyson-scale source; a dead one grinds itself
to dust on the collisional time $P_{\rm orb}/f$ and is gone from any survey
within $P_{\rm orb}/f_{\min}$ of the detection floor---\MaxVisibleResidueMilli
at 2.7\,AU for $f_{\min}=10^{-3}$ \citep{Lacki2025,GooPaper}.  A catastrophe
confined to a planet's surface changes nothing a survey measures.  Visibility
is therefore not a property of how bad the catastrophe was; it is a property
of where it happened and whether it is still happening.

\section{The classes}

Table~\ref{tab:classes} applies the three properties.  We take the classes in
order of the size of their unit, because size is what dispersal depends on.

\begin{table}
\centering
\small
\caption{Self-replicating technologies by the three properties.  Doubling
times are the fastest documented or credibly projected; reach is the extent
an escaped population can attain given its dispersal; visibility is what an
outside astronomer could see.}
\label{tab:classes}
\begin{tabular}{p{2.6cm}p{1.6cm}p{2.5cm}p{2.3cm}p{2.0cm}p{3.0cm}}
\toprule
class & doubling & limit & dispersal & reach & astronomical visibility \\
\midrule
network worm \citep{Moore2003} & 8.5\,s & hosts, bandwidth & the network & the civilisation's hardware & none \\
self-replicating AI agent \citep{Pan2024} & minutes--hours & compute, access & the network; then the economy & hardware, then whatever the economy reaches & none unless it captures physical industry \\
engineered pathogen \citep{Millett2017,Esvelt2022} & hours & host population & hosts' movement & host species; planet & loss of a biosignature \\
mirror life \citep{Adamala2024} & hours & environmental chemistry & the biosphere & planet & loss of a biosignature \\
gene drive \citep{Esvelt2014,Noble2018} & generations & target species & the species' range & species & none \\
molecular assembler \citep{Drexler1986,Freitas2000} & $10^{2}$--$10^{4}$\,s & heat (in place); starlight; feedstock & none / space transport / blowout / ship & planet / system / annulus / galaxy & none / waste heat / seeded systems / front \\
self-reproducing machine, space industry \citep{Freitas1980b,Zykov2005,Metzger2013,Moses2020} & months--years & feedstock, transport & space transport & planetary system & waste heat at the feedstock's orbit \\
von Neumann probe \citep{Freitas1980,Tipler1980} & years & error catastrophe; purpose & starship & galaxy; cosmos & front; Kardashev-III reprocessing \\
\bottomrule
\end{tabular}
\end{table}

\paragraph{Computational replicators.}  Worms have the fastest replication
ever measured in a technological system---the Slammer worm doubled every
8.5\,s and reached $75\,000$ hosts in ten minutes \citep{Moore2003}; a
``flash worm'' with a precomputed target list could do the whole vulnerable
Internet in tens of seconds \citep{Staniford2002}---and the smallest reach:
the substrate is the network, and the network ends at the civilisation's own
hardware.  A self-replicating AI agent \citep{Pan2024} is the same class with
a substrate that includes whatever the network controls; its reach becomes
physical only when it captures physical industry, at which point it is one of
the rows below and its dispersal is theirs.  Astronomically these are
invisible; their catastrophe, if it comes, is the civilisation's, and the sky
does not record it.

\paragraph{Biological replicators.}  Engineered pathogens, mirror organisms
and gene drives are planet-bound by construction: their dispersal is their
hosts' or the biosphere's, their reach the species or the planet, and their
limit the ecology they invade.  Mirror life is the case closest to goo in
spirit---a replicator that the existing biosphere can neither eat nor
resist---and its recent assessment \citep{Adamala2024} reads like
\citet{Freitas2000} with the chemistry changed.  Astronomically the worst case
is the loss of a biosignature, which is a change in one planet's spectrum
that no present survey would notice and that a future one would read as a
planet that never had life.

\paragraph{Molecular assemblers.}  The class with the widest range of
dispersal, and the subject of the companion paper.  Loose on a planet it is
planet-bound, and the planet's own thermodynamics limit it to \CrustGovTime
for the crust.  Loose in space it is system-scale within a century and a
Dyson-scale waste-heat source at the temperature of its feedstock's orbit
(\BeltTempK\,K for an asteroid belt).  If its units are below
\BlowoutRadiusUm\um they leave on starlight; if they inherit a starship they
fill the Galaxy.  The class's accident probability is elicited at
$\sim10^{-3}$--$10^{-4}$ per century \citep{Sandberg2008}, and the companion
paper's central result is that the astronomical record moves that number by
a factor \BranchRatioTypical if one per cent of such accidents would be
interstellar, because the record constrains only the escaping branch.

\paragraph{Macroscopic machines and space industry.}  The self-replicating
lunar factory of \citet{Freitas1980b}, the demonstration replicators of
\citet{Zykov2005}, the bootstrapped space industry of \citet{Metzger2013} and
the review of \citet{Moses2020} describe replicators with doubling times of
months to years and units far too large for passive escape.  Their reach is
the planetary system; their accident is a swarm that eats the small bodies
and, as \citet{Ellery2022b} argues, must be curbed genetically; their
visibility is exactly the waste-heat signature the Dyson searches have looked
for around $5\times10^{6}$ stars \citep{Suazo2024}.  This is the one accidental
class the sky already bounds, and it bounds it while the machines run: once
they stop, the residue is invisible within $P_{\rm orb}/f_{\min}$.

\paragraph{Von Neumann probes.}  The only class with galactic reach, the only
class a civilisation must build on purpose to travel, and the only class the
$10^{5}$-galaxy record constrains: fewer than $\sim\LbExtPersistLamHundred$ of
civilisations can have released a persistent galaxy-spanning replicator if
$\sim100$ civilisations have arisen per galaxy \citep{GooPaper}.  It is also
the class most exposed to the error catastrophe---$10^{4}$ hops across the
Galaxy are $10^{4}$ generations at least---and the one for which
\citet{SaganNewman1983} argued, before nanotechnology had a name, that no
prudent maker would build it.

\section{What follows}

\paragraph{The dispersal principle.}  An accidental replicator inherits the
dispersal of the environment it escapes into, and for a civilisation at our
stage that environment is a planet.  Every class in Table~\ref{tab:classes}
except the last is bounded by the planet or the system, and the last is not
an accident but a project.  The corollary for prioritisation is that the
classes should be ranked by the product of their probability and their
planet-scale consequence, not by their theoretical reach; the reach that
frightens---the Galaxy consumed---belongs to the class that is furthest from
us and easiest to design safely, because we would have to design it at all.

\paragraph{Containment means different things at different scales.}  On a
planet, containment is thermal and chemical: a molecular replicator cannot
process matter faster than the surface sheds heat, and a biological one
cannot outrun the ecology's predators forever; the defender's instruments are
the ones \citet{Freitas2000} proposed, thermal monitoring and quarantine.  In
a system, containment is transport: a swarm that cannot launch cannot spread,
and the design of replication limiters \citep{Ellery2022b} is the whole
game.  Between stars, containment is the error catastrophe, which nature
imposes for free on any copier that is not engineered past it.

\paragraph{The epistemic position.}  A species estimating the probability
that it will die of its own replicators is estimating a risk whose
realisation, in every class but one, leaves no record anyone else could
read.  It cannot learn from the sky, because the sky only records the class
it would have to build on purpose.  It cannot learn from its own past,
because the anthropic shadow \citep{Cirkovic2010} means that a species which
had died this way would not be here to count itself.  It can learn only from
the engineering: from what its replicators can actually do, and from what
stops them.  The companion paper's astronomical bounds are therefore not a
comfort about the planetary classes; they are a demonstration of where
comfort cannot come from.  The one positive statement the record supports is
that a replicator population that keeps growing across a galaxy is not a
common end of civilisations, so that whatever ends them, if anything does, is
something that stays home.

\paragraph{What would change this.}  A detected exoplanet with a refined,
chemically dead surface would put the planet-bound molecular class into the
observational record for the first time.  A waste-heat source whose
temperature matches its host's belt rather than its habitable zone would be
the first system-scale accident seen from outside.  And a demonstrated
replicator of any class with an unbounded substrate---a mirror organism that
grows in the wild, an assembler that eats regolith---would move the problem
from elicitation to measurement, which is where every other risk in the
physical sciences has eventually gone.

\section*{Data and code availability}
The numbers shared with the companion paper are generated by the
\texttt{seti.goo} package of the accompanying repository.

\bibliography{../goo/refs,refs}
\end{document}
